% Unregistered candidate, notes only during the checkpoint14 freeze.
\subsection{Odd abelian factors over two-power dihedral covers}

\paragraph{Candidate conclusion.}
Assuming the separately written and audited A.70, both all-prime
assertions of Task A.5 hold for every actual cover whose geometric
Galois closure is a quotient of
\(D_{2m}\times B\), where \(m\) is a power of two and \(B\)
is an abelian group of odd order. No assertion is made for a
product with a second nonabelian Sylow factor.

\begin{proof}
\paragraph{The actual group and function-field factorization.}
A quotient of a direct product of groups of coprime orders is
the product of the two corresponding quotients. Indeed every
subgroup of such a product splits into its two intersections:
for an element \((x,y)\), suitable powers, using Bezout for the
two group orders, give \((x,1)\) and \((1,y)\) in the subgroup.
Apply this also to the normal kernel of the quotient. Every
quotient of a dihedral two-group is dihedral or abelian; the
abelian cases will be genus-one covers. Thus the nonabelian
case has actual closure group \(G_2\times B_0\), with \(G_2\)
dihedral of two-power order and \(B_0\) abelian of odd order.
Every actual covering subgroup has the form \(H_2\times H_B\).

The abelian cover of \(U\) corresponding to \(B_0/H_B\)
extends to an actual elliptic isogeny
\[
                         \eta:E'\longrightarrow E.
\]
The peripheral element is a commutator and is killed by every
abelian quotient, so this extension is unramified on the proper
elliptic curve. After a finite number-field extension choose an
origin above \(O\); the finite morphism is then an isogeny.
Let \(Z\to E\) be the compactification of the actual
\(H_2\)-intermediate of the dihedral closure. The given cover
has function field equal to the compositum of those of \(Z\)
and \(E'\) over \(E\). This is a field: the normal closure
of the first extension has two-power degree, and the second is
Galois of odd degree, so their intersection is the base field.
Consequently the given compact curve is exactly the connected
normalization of
\[
                         Z\times_E E'.
\]
The removed set is the full inverse image of \(O\), namely the
pairs of a cusp of \(Z\) and a point of \(\eta^{-1}(O)\).
This identifies the actual covering map, not merely its degree
or geometric representation.

\paragraph{Good models of the odd isogeny and of the actual map.}
If \(Z\) has genus one, so does its unramified isogeny pullback,
and A.40 applies. Otherwise A.70, through A.44 or A.69 and the
actual norm maps, gives a potential-good Jacobian at two and a
stable marked model \(\mathcal Z/R\) of compact type in which
every cusp has a distinct companion joined by rational
components. Enlarge the local field at two so that this model,
all marks and the good elliptic model \(\mathcal E/R\) are
defined. The isogenous elliptic curve \(E'\) also has good
reduction, and \(\eta\) extends to an isogeny
\(\mathcal E'\to\mathcal E\) of odd degree. Its finite flat
kernel has order invertible in the residue field, so this
extended isogeny is finite \'{e}tale.

The actual generic map \(q:Z\to E\) extends over
\(\mathcal Z\). Here is a construction avoiding an assertion
about an arbitrary rational map on a singular surface. Use a
marked section \(P\) with \(q(P)=O\). The compact-type
Abel map based at \(P\) extends to
\(\mathcal Z\to\operatorname{Pic}^0(\mathcal Z/R)\), as in
the precise Abel-map input of A.41
\cite[Theorem 1.1 and Proposition 3.2]{fghDupuyRabinoff}.
The actual homomorphism
\(J_Z\to E\) extends between the good abelian schemes.
Their composition restricts generically to \(q\), and is the
required extension. Inserted rational components are mapped
constantly by this same construction.

It follows that
\[
              \mathcal Y=\mathcal Z\times_{\mathcal E}\mathcal E'
\]
is finite \'{e}tale over \(\mathcal Z\). It is a proper flat
semistable model of the actual connected generic pullback.
The generic-fibre field comparison above is what proves its
geometric connectedness; no disconnected formal fibre product
is selected without justification.

\paragraph{Why the pulled-back dual graph is a tree.}
One must not infer compact type just from the finite \'{e}tale
map \(\mathcal Y\to\mathcal Z\): covers of component curves
can join several copies of the original graph and create cycles.
Instead take the actual geometric Galois closure over \(E\).
Its group \(G_2\times B_0\) is nilpotent. Its sole possible
branch point is \(O\), since the odd isogeny is unramified.
Over the good elliptic model the branch closure is an \'{e}tale
section. Generic characteristic zero supplies tameness.
Lehr--Matignon Theorem 6.13, with exactly the scope checked in
A.69, gives a potential-good Jacobian of this actual closure
\cite{fghLehrMatignon2005}.
The actual norm to the given intermediate then gives potential
good reduction of \(J_Y\). After a further finite extension,
the semistable special fibre of \(\mathcal Y\) therefore has a
tree dual graph. No claim about its special two-rank is needed.

\paragraph{Every selected cusp has a lifted rational companion.}
Let a chosen cusp of the pullback lie above a cusp \(P\) of
\(Z\). Choose a distinct companion \(Q\) on \(Z\), and let
\(\mathcal T\) be the connected rational subtree in the
geometric special fibre joining their marked vertices. It
includes both endpoint components. Each component is a
projective line, so its map to the good elliptic curve is
constant; the constants agree across the nodes. Since the
cusp sections map to \(O\), the whole subtree maps to the
single special point \(O\).

The finite \'{e}tale isogeny has, over this algebraically closed
residue field, a discrete fibre of \(\deg\eta\) distinct
points above \(O\). Therefore its pullback over \(\mathcal T\)
is literally the disjoint union of \(\deg\eta\) copies of
\(\mathcal T\). The selected marked section lies on exactly
one such copy. That same copy contains a marked section above
\(Q\): sections above a cusp are the pairs with the kernel
sections of \(\eta\), and prime-to-two kernel sections have
distinct reductions. These two generic marked points are
distinct because their images \(P,Q\) are distinct.

Thus the selected cusp is joined to a genuine companion by an
entirely rational path in \(\mathcal Y\). Stabilization may
contract rational components but cannot insert a positive-genus
component into this path. Since the whole graph is a tree, this
is the unique joining path required by A.42. The argument works
for every selected cusp and does not require the new Jacobian
to have two-rank zero.

A.42 now proves the full fixed-inertia assertion at every
coefficient prime. A.59 proves every normalizer assertion for
the nilpotent geometric closure. As in A.70, isomorphisms,
subgroups and marked points need only a finite field of
definition; geometric base-point paths are used up to inner
conjugacy, and the proved finite-field-restriction argument
supplies every original permitted packet and open derived
subgroup. This completes the candidate deduction.
\end{proof}
